%%%%%%%%%%%%%%%%%%%%%%%%%%%%%%%%%%%%%%%%%%%%%%%%%%%%%%%%%%%%%%%%%%%%%%%%%%%%%%
% Short Sectioned Assignment LaTeX Template Version 1.0 (5/5/12)
% This template has been downloaded from: http://www.LaTeXTemplates.com
% Original author:  Frits Wenneker (http://www.howtotex.com)
% License: CC BY-NC-SA 3.0 (http://creativecommons.org/licenses/by-nc-sa/3.0/)
%%%%%%%%%%%%%%%%%%%%%%%%%%%%%%%%%%%%%%%%%%%%%%%%%%%%%%%%%%%%%%%%%%%%%%%%%%%%%%

% \documentclass[paper=a4, fontsize=11pt]{scrartcl} % A4 paper and 11pt font size
\documentclass[11pt, a4paper]{book}
\usepackage[T1]{fontenc} % Use 8-bit encoding that has 256 glyphs
\usepackage[utf8]{inputenc}
\usepackage{fourier} % Use the Adobe Utopia font for the document - comment this line to return to the LaTeX default
\usepackage[spanish, es-tabla]{babel} % Selecciona el español para palabras introducidas automáticamente, p.ej. "septiembre" en la fecha y especifica que se use la palabra Tabla en vez de Cuadro
\usepackage{url} % ,href} %para incluir URLs e hipervínculos dentro del texto (aunque hay que instalar href)
\usepackage{graphics, graphicx, float} %para incluir imágenes y colocarlas
\usepackage[gen]{eurosym} %para incluir el símbolo del euro
\usepackage{cite} %para incluir citas del archivo <nombre>.bib
\usepackage{enumerate}
\usepackage{hyperref}
\usepackage{graphicx}
\usepackage{tabularx}
\usepackage{booktabs}
\usepackage{amsmath}

\usepackage[table,xcdraw]{xcolor}
\hypersetup{
	colorlinks=true,	% false: boxed links; true: colored links
	linkcolor=black,	% color of internal links
	urlcolor=cyan		% color of external links
}

\renewcommand{\familydefault}{\sfdefault}
\usepackage{fancyhdr} % Custom headers and footers
\pagestyle{fancyplain} % Makes all pages in the document conform to the custom headers and footers
\fancyhead[L]{} % Empty left header
\fancyhead[C]{} % Empty center header
\fancyhead[R]{Francisco David Castejón Soto} % My name
\fancyfoot[L]{} % Empty left footer
\fancyfoot[C]{} % Empty center footer
\fancyfoot[R]{\thepage} % Page numbering for right footer
\renewcommand{\headrulewidth}{0pt} % Remove header underlines
\renewcommand{\footrulewidth}{0pt} % Remove footer underlines
\setlength{\headheight}{13.6pt} % Customize the height of the header

\usepackage{titlesec, blindtext, color}
\definecolor{gray75}{gray}{0.75}
\newcommand{\hsp}{\hspace{20pt}}
\titleformat{\chapter}[hang]{\Huge\bfseries}{\thechapter\hsp\textcolor{gray75}{|}\hsp}{0pt}{\Huge\bfseries}
\setcounter{secnumdepth}{4}
\usepackage[Lenny]{fncychap}

% For code snippets
\usepackage{minted}
\usemintedstyle{colorful}
\usepackage{caption}

\begin{document}

% Páginas preliminares sin numeración
\frontmatter

\begin{titlepage}
	\thispagestyle{empty}

	\centering
	\includegraphics[width=0.9\textwidth]{logos/logo_ugr.jpg}
	\vspace{1.0cm}

	{\rmfamily\textsc{\Large MÁSTER UNIVERSITARIO}}\\[0.2cm]
	{\rmfamily\textsc{\large en Ciencia de Datos e Ingeniería de Computadores}}
	\vspace{1.5cm}

	\vfill

	{\Huge\bfseries Trabajo de Fin de Máster}\\[0.5cm]
	\rule{\textwidth}{3pt}\\[0.5cm]
	{\Large\bfseries Desarrollo y Optimización de un Agente Inteligente \\
	para Scripts of Tribute mediante Algoritmos Evolutivos}
	\vspace{1.5cm}

	\vfill

	\begin{tabular}{@{}c@{}}
		\textbf{\large Autor}         \\[0.3cm]
		Francisco David Castejón Soto \\[1cm]
		\textbf{\large Director}      \\[0.3cm]
		Dr. Pablo García Sánchez
	\end{tabular}

	\vfill

	\includegraphics[width=0.4\textwidth]{logos/etsiit_logo.png}\\[0.3cm]
	{\rmfamily\textsc{\footnotesize Escuela Técnica Superior de Ingenierías Informática y de Telecomunicación}}\\
	\rule{0.1\textwidth}{0.5pt}\\[0.3cm]
	{\large Granada, 1 Julio de 2025}

\end{titlepage}
\chapter*{Prefacio}

Para la realización de este documento, se ha utilizado la plantilla \LaTeX \cite{guervos_jjplantilla-tfg-etsiit_2024} específica para la ETSIIT.

% Índices
\cleardoublepage
\tableofcontents

\cleardoublepage
\listoffigures

\cleardoublepage
\listoftables

% Contenido principal
\mainmatter

\chapter{Introducción}


\section{Antecedentes}


\section{Motivación}


\section{Objetivos}


\section{Estructura de la memoria}



\chapter{Estado del arte}


\section{Dominio del problema}


\section{Trabajos relacionados}



\chapter{Propuesta}


\section{Metodología de desarrollo}


\section{Planificación temporal}


\section{Presupuesto}
\chapter{Desarrollo}

\section{Sección 1...}

\chapter{Experimentos realizados}

\section{Sección 1...}

\chapter{Resultados}

\section{Sección 1...}

\chapter{Conclusiones}

\section{Trabajos futuros}


% Material de referencia
\backmatter

% % Bibliografía
\cleardoublepage
\bibliographystyle{plainurl}
\bibliography{bibliografia}

% \appendix
% \input{apendices/codigo}

\end{document}